\section{Introducción: por qué se necesita RL para el razonamiento}

Esta clase la imparte como profesor invitado Aviral Kumar, de CMU, y se centra en la aplicación del aprendizaje por refuerzo (RL) al razonamiento en los modelos de lenguaje de gran tamaño (LLM). El ponente acota el razonamiento al escenario de resolución de problemas matemáticos y organiza el contenido en dos grandes partes a lo largo de una línea de tiempo: las técnicas clásicas de RL y las extensiones modernas.

\begin{importantbox}{Motivación central}
El paradigma tradicional de entrenamiento basado en next-token prediction (NTP) tiene una limitación fundamental: el error entre el modelo aprendido $\hat{p}_\theta(\mathbf{y}|\mathbf{x})$ y la distribución real $p^*(\mathbf{y}|\mathbf{x})$ es proporcional a $|\mathcal{D}(\mathbf{y}|\mathbf{x})|^{-\alpha}$; es decir, el error disminuye conforme aumenta la cantidad de datos similares a la entrada objetivo $\mathbf{x}$. En los problemas de razonamiento, los datos de alta calidad de pares problema--solución son sumamente escasos, y se prevé que hacia 2028 se agoten los datos de texto de alta calidad disponibles en internet.
\end{importantbox}

\subsection{Modos de falla de NTP}

Cuando un modelo entrenado solo de forma supervisada resuelve problemas matemáticos difíciles, como los del AIME (American Invitational Mathematics Examination) o la IMO (Olimpiada Internacional de Matemáticas), puede generar texto que «parece una solución», pero su cadena de razonamiento lógico suele contener errores críticos. El modelo omite pasos de razonamiento necesarios o, durante una demostración, analiza términos aislados y pasa por alto la estructura global.

\begin{figure}[H]
\centering
\includegraphics[width=0.9\textwidth]{slides-images/slide-03.jpg}
\caption{Entrenamiento tradicional: la cota de error de next-token prediction\protect\footnotemark}
\end{figure}
\footnotetext{Slides, página 3.}

\begin{figure}[H]
\centering
\includegraphics[width=0.9\textwidth]{slides-images/slide-04.jpg}
\caption{El problema de la escasez de datos: limitaciones de datos en el razonamiento matemático y en la IA corporizada\protect\footnotemark}
\end{figure}
\footnotetext{Slides, página 4.}

\subsection{Resumen de la sección}

La razón principal por la que NTP falla en escenarios de razonamiento es la insuficiencia de datos de razonamiento de alta calidad. El valor de RL consiste en que no depende de soluciones completas anotadas por personas, sino que guía al modelo mediante una señal de recompensa (por ejemplo, si la respuesta es correcta) para que descubra por sí mismo las rutas de razonamiento correctas.
